\documentclass[11pt,a4paper]{article}

\usepackage[utf8]{inputenc}
\usepackage[T1]{fontenc}
\usepackage[provide=*,french]{babel}
\usepackage[margin=2.3cm]{geometry}
\usepackage{titlesec}
\usepackage{enumitem}
\usepackage{array}
\usepackage{longtable}
\usepackage{booktabs}
\usepackage{fancyhdr}
\usepackage{hyperref}
\usepackage{tabularx}
\usepackage{amsmath}

\hypersetup{
  colorlinks=true,
  linkcolor=black,
  urlcolor=black,
  citecolor=black,
  pdftitle={Protocole de collecte de donnees v1},
  pdfauthor={}
}

% --- Mise en forme sobre, noir et blanc, lignes fines ---
\titleformat{\section}
  {\normalfont\Large\bfseries}
  {\thesection}{0.8em}{}
\titlespacing{\section}{0pt}{1.4em}{0.8em}

\titleformat{\subsection}
  {\normalfont\large\bfseries}
  {\thesubsection}{0.7em}{}
\titlespacing{\subsection}{0pt}{1.1em}{0.5em}

\titleformat{\subsubsection}
  {\normalfont\bfseries\itshape}
  {\thesubsubsection}{0.6em}{}
\titlespacing{\subsubsection}{0pt}{0.9em}{0.4em}

\pagestyle{fancy}
\fancyhf{}
\fancyfoot[C]{\small Protocole de collecte v1 --- rPPG / rBCG}
\fancyfoot[R]{\small\thepage}
\renewcommand{\headrulewidth}{0pt}
\renewcommand{\footrulewidth}{0.3pt}

\setlist[itemize]{leftmargin=*,itemsep=2pt,topsep=3pt}
\setlist[enumerate]{leftmargin=*,itemsep=2pt,topsep=3pt}

% --- Encadres sobres : simple cadre noir, fond blanc ou gris tres clair ---
\usepackage[most]{tcolorbox}

\newtcolorbox{notebox}[1]{
  colback=white, colframe=black, boxrule=0.6pt, arc=0pt,
  left=8pt, right=8pt, top=5pt, bottom=5pt,
  title={#1}, coltitle=black, fonttitle=\bfseries\small,
  attach boxed title to top left={yshift=-2mm, xshift=2mm},
  boxed title style={colback=white, boxrule=0pt, top=1pt, bottom=1pt},
  breakable
}

\newtcolorbox{warnbox}[1]{
  colback=white, colframe=black, boxrule=1pt, arc=0pt,
  left=8pt, right=8pt, top=5pt, bottom=5pt,
  title={#1}, coltitle=black, fonttitle=\bfseries\small,
  attach boxed title to top left={yshift=-2mm, xshift=2mm},
  boxed title style={colback=white, boxrule=0pt, top=1pt, bottom=1pt},
  breakable
}

\newtcolorbox{codebox}{
  colback=black!4, colframe=black!4, boxrule=0pt, arc=0pt,
  left=8pt, right=8pt, top=4pt, bottom=4pt,
  fontupper=\ttfamily\small, breakable
}

\newcolumntype{P}[1]{>{\raggedright\arraybackslash}p{#1}}

\begin{document}

% ================================================================
% PAGE DE GARDE
% ================================================================
\begin{center}
{\LARGE\bfseries PROTOCOLE DE COLLECTE DE DONN\'EES}\\[8pt]
{\Large Version 1 --- D\'etaill\'e}\\[6pt]
{\large Screening cardiovasculaire non invasif par smartphone (rPPG / rBCG)}\\[4pt]
{\normalsize R\'eseau de pharmacies --- Afrique subsaharienne}
\end{center}

\vspace{0.6cm}
\noindent\rule{\textwidth}{0.6pt}
\vspace{0.3cm}

\begin{notebox}{Pourquoi ce protocole}
Les m\'ethodes de mesure \`a distance du rythme cardiaque (rPPG) ont \'et\'e
d\'evelopp\'ees et valid\'ees presque exclusivement sur des populations \`a
peau claire, en conditions de laboratoire. Une analyse r\'ecente de la
litt\'erature (Dasari et al., 2021, repris dans Mehmood et al., 2025) montre
que l'erreur moyenne (MAE) de la m\'ethode classique CHROM passe de
5.2~bpm sur les phototypes Fitzpatrick I--III \`a 14.1~bpm sur les
phototypes V--VI --- une d\'egradation qui rend la mesure cliniquement
inutilisable. Les mod\`eles d'apprentissage profond (PhysNet, TS-CAN) sont
moins affect\'es (6.0 $\rightarrow$ 9.5~bpm) mais restent insuffisants
(section~4).

Ce protocole vise \`a constituer le premier dataset de r\'ef\'erence collect\'e
en conditions r\'eelles de pharmacie africaine, avec une priorit\'e absolue
donn\'ee aux phototypes V--VI, afin de permettre le fine-tuning de mod\`eles
qui fonctionnent r\'eellement pour cette population.
\end{notebox}

\vspace{0.3cm}
\tableofcontents
\newpage
\section{Contexte scientifique et objectifs}

\subsection{Le probl\`eme physique : pourquoi les peaux fonc\'ees posent
probl\`eme}

Le signal rPPG repose sur la d\'etection de variations de couleur de la peau
de l'ordre de 0.1\% de la valeur d'un pixel, caus\'ees par l'afflux sanguin
\`a chaque battement cardiaque. Ce signal est extr\^emement faible et doit
\^etre extrait d'un bruit 10 \`a 500 fois plus important (variations
d'\'eclairage, r\'eflexion sp\'eculaire, mouvement).

La m\'elanine, pr\'esente en plus grande quantit\'e dans les phototypes
Fitzpatrick \'elev\'es, absorbe une partie significative de la lumi\`ere qui
transporte ce signal. Concr\`etement~:

\begin{itemize}
\item Le rapport signal/bruit (SNR) de la composante diffuse --- celle qui
porte l'information cardiaque --- diminue d'environ 9.5~dB pour un phototype~I
\`a environ 4.5~dB pour un phototype~VI.
\item Sur smartphone bas de gamme, avec une dynamique de capteur limit\'ee,
les pixels d'une peau tr\`es fonc\'ee peuvent \^etre satur\'es en valeurs
basses (proches de 0) sous \'eclairage faible, ou en valeurs hautes
(proches de 255) sous \'eclairage trop fort. Dans les deux cas, le signal
de 0.1\% devient invisible car le pixel ne varie plus du tout.
\end{itemize}

\begin{warnbox}{Cons\'equence directe pour le protocole}
Pour les phototypes V et VI, l'\'eclairage doit \^etre maintenu dans une
plage interm\'ediaire --- ni trop faible (saturation basse), ni trop fort
(saturation haute). La section~5.4 d\'etaille comment v\'erifier cela en
temps r\'eel pendant la collecte.
\end{warnbox}

\subsection{Le probl\`eme algorithmique : pourquoi CHROM \'echoue}

La m\'ethode CHROM (De Haan \& Jeanne, 2013) repose sur l'hypoth\`ese que
toute peau humaine, sous lumi\`ere blanche, poss\`ede approximativement le
m\^eme vecteur de couleur normalis\'e~:
\[
\left[\frac{R}{\|\cdot\|}, \frac{G}{\|\cdot\|}, \frac{B}{\|\cdot\|}\right]
\approx [0.7682,\ 0.5121,\ 0.3841]
\]
o\`u $\|\cdot\| = \sqrt{R^2+G^2+B^2}$. Les coefficients $3, 2, 1.5, 1, 1.5$ de
la formule CHROM sont directement d\'eriv\'es de ce vecteur (voir
section~7.2 pour la d\'erivation compl\`ete).

Ce vecteur a \'et\'e mesur\'e sur des sujets de phototypes~I \`a~V dans un
laboratoire europ\'een, sous \'eclairage blanc contr\^ol\'e. Deux \'ecarts
se cumulent en conditions r\'eelles africaines~:

\begin{enumerate}
\item \textbf{\'Ecart li\'e au phototype} : la m\'elanine absorbe
pr\'ef\'erentiellement le vert et le bleu. Le vrai vecteur de peau pour un
phototype~VI est probablement plus proche de $[0.85,\ 0.43,\ 0.30]$ que de
$[0.7682,\ 0.5121,\ 0.3841]$.
\item \textbf{\'Ecart li\'e \`a l'\'eclairage} : le vecteur standard suppose
une lumi\`ere blanche. Sous LED chaude ou n\'eon, le spectre de la lumi\`ere
ambiante d\'eforme diff\'eremment les trois canaux R, G, B, ce qui modifie
le ratio optimal entre les deux signaux de chrominance.
\end{enumerate}

Quand les coefficients de la formule ne correspondent plus au vrai vecteur
de peau et \`a la vraie couleur de la lumi\`ere, la projection de chrominance
n'annule plus correctement le bruit lumineux --- un r\'esidu de lumi\`ere
reste dans le signal final et peut devenir plus grand que le signal cardiaque
lui-m\^eme. C'est ce r\'esidu que ce protocole vise \`a mesurer et \`a
corriger par phototype.

\subsection{Objectifs du dataset}

\begin{enumerate}
\item Mesurer objectivement, pour chaque phototype, l'\'ecart entre le
vecteur de peau r\'eel et le vecteur standard de De Haan, afin de recalculer
des coefficients CHROM adapt\'es (section~7).
\item Fournir des donn\'ees d'entra\^inement et de fine-tuning pour les
mod\`eles PhysNet et TS-CAN, dans des conditions repr\'esentatives du
d\'eploiement final.
\item Couvrir la diversit\'e r\'eelle des conditions de d\'eploiement~:
\'eclairage de pharmacie (LED, n\'eon), lumi\`ere naturelle tropicale,
smartphones d'entr\'ee de gamme, et niveaux de mouvement/parole r\'ealistes.
\item Permettre une \'evaluation rigoureuse, strat\'ifi\'ee par phototype et
par condition, des performances de chaque m\'ethode (CHROM original, CHROM
adaptatif, PhysNet, TS-CAN).
\end{enumerate}



\section{Mesures collect\'ees}

\subsection{D\'efinitions}

\subsubsection*{PPG --- Photoplethysmogram}

Le PPG est la forme d'onde brute du pouls, telle que mesur\'ee par
l'oxym\`etre. C'est une courbe continue qui oscille \`a chaque battement
cardiaque. Le PPG est la donn\'ee la plus riche~: HR, HRV et SpO2 en sont
tous d\'eriv\'es. Pour l'entra\^inement des mod\`eles par apprentissage
profond (loss de corr\'elation de Pearson, section~7.3), disposer de la
forme d'onde compl\`ete et pas seulement d'un HR ponctuel est
essentiel.

\subsubsection*{HR --- Heart Rate}

Nombre de battements cardiaques par minute (bpm), calcul\'e \`a partir des
pics du PPG. C'est la m\'etrique de sortie principale de toutes les
m\'ethodes \'etudi\'ees (CHROM, POS, PhysNet, TS-CAN).

\subsubsection*{SpO2 --- Saturation en oxyg\`ene}

Pourcentage d'h\'emoglobine satur\'ee en oxyg\`ene dans le sang art\'eriel
(normal~: 95--100\%). Mesur\'e par l'oxym\`etre via deux longueurs d'onde
(rouge et infrarouge). Pertinent pour le d\'epistage de pathologies
respiratoires.

\subsubsection*{BP --- Blood Pressure}

Pression art\'erielle, exprim\'ee en mmHg sous la forme
systolique/diastolique (ex.\ 120/80). Mesur\'ee par un tensiom\`etre \`a
brassard. C'est la cible de l'estimation par Pulse Transit Time (PTT,
section~8).

\subsection{Tableau r\'ecapitulatif}

\begin{center}
\begin{tabularx}{\textwidth}{P{2.2cm} X P{4.3cm}}
\toprule
\textbf{Mesure} & \textbf{Description} & \textbf{Mat\'eriel} \\
\midrule
PPG & Forme d'onde compl\`ete du pouls, \'echantillonn\'ee \`a $\sim$60~Hz &
Oxym\`etre de pouls (type Contec CMS50D+, USB) \\
HR & Fr\'equence cardiaque instantan\'ee, \'echantillonn\'ee \`a 1~Hz &
Oxym\`etre (calcul interne) \\
SpO2 & Saturation en oxyg\`ene, \'echantillonn\'ee \`a 1~Hz &
Oxym\`etre (calcul interne) \\
BP & Pression art\'erielle syst./diast., une mesure par enregistrement~2 &
Tensiom\`etre \`a brassard (type Omron) \\
Vid\'eo visage & Signal source pour rPPG (couleur) et rBCG (mouvement) &
Smartphone, cam\'era avant, 30--60~fps \\
Vid\'eo paume (optionnel) & Second point de mesure pour le PTT multi-site &
Smartphone, cam\'era avant \\
\bottomrule
\end{tabularx}
\end{center}

\subsection{Pourquoi ces quatre mesures, ni plus ni moins}

Le choix de limiter ce protocole \`a PPG, HR, SpO2 et BP en excluant
ECG, fr\'equence respiratoire d\'edi\'ee et r\'eponse \'electrodermale
suit directement la logique du protocole VitalVideos (Toye, 2023), qui est
notre principale source d'inspiration m\'ethodologique pour cette section
(voir section~3 pour le d\'etail des emprunts).

\begin{itemize}
\item \textbf{PPG, HR et SpO2 proviennent du m\^eme appareil, au m\^eme
co\^ut.} L'oxym\`etre qui capture le PPG calcule le HR et le SpO2 en
interne, \`a partir des deux longueurs d'onde (rouge/infrarouge) qu'il
utilise d\'ej\`a. Ne pas les enregistrer reviendrait \`a ignorer une
donn\'ee qui n'a aucun co\^ut marginal de collecte.
\item \textbf{BP est un signal physiologiquement ind\'ependant du PPG
seul.} Le PPG capture la forme du flux sanguin en un point. La pression
art\'erielle d\'epend de la rigidit\'e art\'erielle et du volume sanguin
total une information que le PPG \`a un point unique ne contient pas.
VitalVideos a donc ajout\'e un tensiom\`etre s\'epar\'e plut\^ot que de
tenter de d\'eriver le BP du PPG seul, pr\'ecis\'ement parce que ce sont
deux mesures ind\'ependantes qui doivent \^etre collect\'ees ensemble pour
qu'un mod\`ele apprenne leur relation (typiquement via le PTT, qui
n\'ecessite par d\'efinition deux points de mesure distincts, pas un seul
PPG).
\item \textbf{ECG, RR d\'edi\'ee et GSR amélioreraient la richesse
scientifique mais au prix d'un mat\'eriel et d'un temps suppl\'ementaires
par sujet.} VitalVideos identifie lui-m\^eme l'ECG comme l'ajout le plus
pr\'ecieux qu'il regrette de ne pas avoir collect\'e (Toye, 2023, section
\guillemotleft~Reflection and suggestions for future datasets~\guillemotright),
suivi de la GSR puis de la RR. Ce protocole fait le choix inverse~:
maximiser le nombre de sujets collect\'es --- en particulier dans les
phototypes V--VI sous-repr\'esent\'es --- plut\^ot que la richesse de
signaux par sujet. PPG/HR/SpO2/BP reste donc le sous-ensemble minimal
cohérent~: retirer SpO2 ou HR ne simplifierait rien (m\^eme geste, m\^eme
appareil), et retirer BP supprimerait l'objectif PTT dans son ensemble.
\end{itemize}

\newpage
\section{Niveaux de collecte}

\subsection{Document d'inspiration et r\'ef\'erences principales}

Ce protocole s'appuie sur trois sources principales, dont les apports
respectifs sont d\'etaill\'es ci-dessous.

\textbf{Toye, P. (2023). \emph{Vital Videos: A public dataset of videos
with PPG and BP ground truths}.} Ce document d\'ecrit le protocole de
collecte du dataset VitalVideos-Europe-1 (850~sujets) et constitue la
principale source d'inspiration m\'ethodologique de ce protocole. Plusieurs
\'el\'ements en sont repris directement~: la s\'eparation en deux
enregistrements distincts (repos puis mesure de pression, section~9.3--9.4),
la r\`egle de placement de l'oxym\`etre sur le bras oppos\'e au brassard
(section~9.2), le seuil de 150~lux pour l'\'eclairage minimal (section~5.4),
et le choix des quatre mesures PPG/HR/SpO2/BP (section~2.3). Les limites que
les auteurs identifient eux-m\^emes --- usage d'une cam\'era industrielle
plut\^ot que d'un smartphone, absence de revue \'ethique formelle,
classification subjective du phototype, sc\'enarios uniquement silencieux,
population europ\'eenne \`a 90\% phototypes~I--III --- sont autant de points
que ce protocole cherche \`a corriger (section~3.6).

\textbf{Dasari, A., Prakash, S. K. A., Jeni, L. A. \& Tucker, C. S. (2021).
\emph{Evaluation of biases in remote photoplethysmography methods}. npj
Digital Medicine.} Cette \'etude est la premi\`ere \`a avoir quantifi\'e
exp\'erimentalement le biais de phototype sur des sujets r\'eels, en
collectant des donn\'ees en Inde et en Sierra Leone et en les comparant \`a
un appareil de r\'ef\'erence approuv\'e FDA. C'est la source des chiffres de
d\'egradation utilis\'es dans ce protocole (section~1, section~4).

\textbf{Mehmood, A. et al. (2025). \emph{Demographic bias in public remote
photoplethysmography datasets}. npj Digital Medicine.} Cette analyse
transversale de 100~\'etudes rPPG synth\'etise les r\'esultats de Dasari et
al.\ et d'autres travaux pour quantifier la sous-repr\'esentation des
phototypes fonc\'es dans les datasets publics existants (section~4).

\subsection{Le dilemme fondamental}

En apprentissage automatique, un principe central veut que la distribution
des donn\'ees d'entra\^inement corresponde \`a la distribution des donn\'ees
de d\'eploiement. Appliqu\'e na\"ivement, ce principe sugg\'ererait de
collecter uniquement en conditions r\'eelles de pharmacie. Mais deux
probl\`emes contradictoires apparaissent~:

\begin{itemize}
\item \textbf{Si la collecte est uniquement en conditions r\'eelles}~: le
patient bouge, parle, l'\'eclairage varie pendant l'enregistrement. Le signal
oxym\`etre de r\'ef\'erence lui-m\^eme devient bruit\'e (artefacts de
mouvement sur le capteur). Le mod\`ele apprend alors une relation
signal/label impr\'ecise --- il apprend en partie le bruit.
\item \textbf{Si la collecte est uniquement en conditions contr\^ol\'ees}~:
les labels sont tr\`es propres, mais le mod\`ele n'a jamais vu la
variabilit\'e r\'eelle pendant l'entra\^inement. C'est exactement le
sch\'ema qui explique l'effondrement de PhysNet (MAE 2.95~bpm sur UBFC,
contr\^ol\'e) lorsqu'il est \'evalu\'e sur MMPD (MAE 4.80~bpm, mobile) ---
et plus encore sur des peaux qu'il n'a jamais vues (MAE 11.82~bpm sur
UCLA-rPPG, peaux mixtes).
\end{itemize}

La solution retenue est une collecte en trois niveaux, combin\'es dans
l'entra\^inement par une strat\'egie progressive (section~9).

\subsection{R\'epartition des trois types}

\begin{center}
\begin{tabularx}{\textwidth}{P{1.6cm} P{1.2cm} X P{1.6cm}}
\toprule
\textbf{Type} & \textbf{Part} & \textbf{Conditions} & \textbf{Sujets} \\
\midrule
A & 20\% & Int\'erieur, \'eclairage LED neutre stable (ring light), sujet
assis immobile, distance fixe, silence & 15--20 \\
B & 50\% & \'Eclairage r\'eel de la pharmacie (LED/n\'eon), sujet assis
naturellement, peut parler doucement & 35--50 \\
C & 30\% & Ext\'erieur ou int\'erieur avec lumi\`ere naturelle, conditions
vari\'ees, l\'egers mouvements naturels & 20--30 \\
\bottomrule
\end{tabularx}
\end{center}

\subsection{TYPE A --- Conditions contr\^ol\'ees (fondation)}

\textbf{Conditions exactes.} Salle ou coin d\'edi\'e de la pharmacie, ring
light LED blanc neutre plac\'e \`a environ 60~cm du visage, sujet assis sur
une chaise fixe, smartphone sur support \`a distance constante (50--70~cm),
regard droit devant, silence pendant l'enregistrement~1.

\textbf{Ce que cela apporte.}
\begin{itemize}
\item Des labels de r\'ef\'erence tr\`es pr\'ecis (variation HR
$\pm$1~bpm), car le signal oxym\`etre n'est pas perturb\'e par le mouvement.
\item Une mesure fiable du vecteur de peau par sujet (section~7), puisque
l'\'eclairage est stable et document\'e.
\item Une base de comparaison directe avec la litt\'erature, les conditions
\'etant proches de celles d'UBFC --- avec la diff\'erence cruciale que les
sujets sont \`a phototypes V--VI.
\end{itemize}

\textbf{Pourquoi seulement 20\%.} Ces donn\'ees apprennent au mod\`ele la
relation fondamentale signal/fr\'equence cardiaque sans ambigu\"it\'e, mais
ne contiennent pas la variabilit\'e du d\'eploiement r\'eel. Une proportion
trop \'elev\'ee risquerait un sur-ajustement \`a des conditions qui ne se
reproduiront jamais en pharmacie.

\subsection{TYPE B --- Conditions semi-r\'eelles (cœur du dataset)}

\textbf{Conditions exactes.} Espace normal de la pharmacie (comptoir, salle
d'attente), \'eclairage tel quel --- pas de ring light additionnel ---,
sujet assis dans sa position naturelle, smartphone tenu par le collecteur ou
pos\'e sur un support simple, parole douce autoris\'ee pendant
l'enregistrement.

\textbf{Sous-types \`a documenter} (pour analyse ult\'erieure par condition
d'\'eclairage)~:
\begin{itemize}
\item B1 (environ 40\% du TYPE B) : \'eclairage n\'eon fluorescent
\item B2 (environ 40\% du TYPE B) : \'eclairage LED chaud
\item B3 (environ 20\% du TYPE B) : m\'elange lumi\`ere naturelle et
artificielle (proximit\'e d'une fen\^etre)
\end{itemize}

\textbf{Ce que cela apporte.} C'est la condition la plus proche du
d\'eploiement final tout en gardant des labels encore fiables --- le sujet
reste coop\'eratif et globalement stable. La variabilit\'e d'\'eclairage
naturelle de la pharmacie (qui change au fil de la journ\'ee, qui peut
inclure des n\'eons d\'efectueux, etc.) est captur\'ee et document\'ee.

\subsection{TYPE C --- Conditions r\'eelles (cible du d\'eploiement)}

\textbf{Conditions exactes.} Smartphone tenu \`a la main (mouvement naturel
de la main inclus), sujet en situation de consultation normale, peut parler,
tourner l\'eg\`erement la t\^ete, \'eclairage non contr\^ol\'e --- ce qui
existe r\'eellement \`a cet endroit et \`a ce moment.

R\'epartition recommand\'ee pour couvrir la diversit\'e de la lumi\`ere
naturelle tropicale (d\'etails en section~6)~:
\begin{itemize}
\item Moments de la journ\'ee~: matin (8h--10h), midi (12h--14h), fin de
journ\'ee (16h--18h)
\item M\'et\'eo~: ensoleill\'e et nuageux
\item Position relative au soleil~: ombre et plein soleil, en privil\'egiant
la lumi\`ere lat\'erale
\end{itemize}

\textbf{Ce que cela apporte.} C'est la distribution exacte du d\'eploiement
r\'eel. Aucun dataset existant --- UBFC, PURE, MMPD, UCLA-rPPG, VitalVideos
--- ne contient cette combinaison sp\'ecifique de population (phototypes
V--VI), de mat\'eriel (smartphones d'entr\'ee de gamme africains) et de
conditions (pharmacie + lumi\`ere tropicale). M\^eme 20--30 sujets bien
document\'es dans ces conditions apportent plus de valeur, pour l'objectif
de d\'eploiement, que des centaines de sujets suppl\'ementaires en
conditions parfaites.

\textbf{Limite \`a documenter.} Les labels de r\'ef\'erence peuvent \^etre
localement moins pr\'ecis (artefacts de mouvement sur l'oxym\`etre). Chaque
enregistrement TYPE C doit \^etre annot\'e avec un indicateur de qualit\'e
du signal de r\'ef\'erence (section~5.5).

\subsection{Ce que ce protocole corrige par rapport \`a VitalVideos}

\begin{center}
\begin{tabularx}{\textwidth}{X X}
\toprule
\textbf{Protocole VitalVideos (Toye, 2023)} & \textbf{Ce protocole} \\
\midrule
Cam\'era industrielle (Basler) sur tr\'epied fixe, mise au point manuelle &
Smartphones r\'eels du march\'e local, repr\'esentatifs du d\'eploiement
final \\
Aucune revue \'ethique formelle (absence d'affiliation institutionnelle) &
Validation \'ethique recherch\'ee via affiliation acad\'emique et/ou
entreprise \\
Classification du phototype par jugement visuel subjectif uniquement &
Classification visuelle compl\'et\'ee par un vecteur de peau mesur\'e
objectivement depuis les pixels (section~4.4) \\
Participants instruits de rester silencieux et immobiles dans tous les cas &
Trois niveaux de mouvement/parole document\'es (TYPE A/B/C) \\
Population europ\'eenne (environ 90\% phototypes I--III) &
Priorit\'e explicite aux phototypes V--VI (section~4) \\
Collecte int\'erieure uniquement (contrainte m\'et\'eorologique) &
Collecte ext\'erieure incluse, avec grille de r\'epartition d\'edi\'ee
(section~6) \\
\bottomrule
\end{tabularx}
\end{center}

\newpage
\section{Distribution des phototypes}

\subsection{Justification de la priorisation : ce que montrent les
\'etudes}

La litt\'erature r\'ecente fournit des chiffres pr\'ecis qui justifient la
priorisation retenue dans ce protocole.

\textbf{Dasari et al.\ (2021, npj Digital Medicine)} ont \'et\'e les
premiers \`a quantifier ce biais directement sur des sujets r\'eels, en
collectant des donn\'ees de fr\'equence de pouls et de vid\'eo faciale auprès
de sujets en Inde et en Sierra Leone, compar\'ees \`a un appareil de
r\'ef\'erence approuv\'e par la FDA. Leurs r\'esultats, repris dans l'analyse
transversale de Mehmood et al.\ (2025, npj Digital Medicine) portant sur
100~\'etudes rPPG, montrent que les m\'ethodes classiques bas\'ees sur la
chrominance (CHROM, POS) ont une erreur moyenne absolue (MAE) de 5.2~bpm sur
les sujets phototypes~I--III, qui se d\'egrade \`a 14.1~bpm sur les sujets
phototypes~V--VI ($p<0.01$). Les mod\`eles d'apprentissage profond montrent
une d\'egradation moindre mais toujours significative, passant de 6.0~bpm \`a
9.5~bpm sur la m\^eme plage de phototypes. Les auteurs concluent que,
m\^eme si les approches neuronales modernes att\'enuent partiellement les
biais li\'es au phototype, les d\'es\'equilibres d\'emographiques des
datasets publics se traduisent toujours par des \'ecarts de pr\'ecision
cliniquement significatifs.

\textbf{Une m\'eta-analyse compl\'ementaire}, portant sur trois datasets
totalisant 73~individus et plus de 400~vid\'eos, rapporte un RMSE passant de
6.69~bpm \`a 18.98~bpm pour les phototypes les plus fonc\'es --- confirmant
que la d\'egradation n'est pas un artefact isol\'e \`a une seule \'etude
mais un r\'esultat reproduit \`a travers plusieurs sources de donn\'ees
ind\'ependantes.

\textbf{La cause structurelle} de ce biais est la sous-repr\'esentation des
phototypes fonc\'es dans les datasets d'entra\^inement existants~: les
phototypes~V et~VI repr\'esentent typiquement moins de 5\% des sujets dans
UBFC, MMSE-HR et AFRL, contre 90\% de phototypes~I--III dans VitalVideos-
Europe-1 (Toye, 2023). Un mod\`ele entra\^in\'e sur ces distributions n'a
statistiquement presque jamais vu de signal rPPG sur peau fonc\'ee pendant
son apprentissage.

\subsection{Pourquoi la d\'egradation n'est pas lin\'eaire}

Un point important pour comprendre la priorisation~: l'erreur n'augmente
pas progressivement du phototype~I au phototype~VI, elle s'acc\'el\`ere
fortement \`a partir du phototype~V. Cela correspond \`a la fois \`a un
effet physique (la concentration de m\'elanine augmente non lin\'eairement
sur l'\'echelle Fitzpatrick) et \`a un effet de distribution des donn\'ees
(les phototypes~V--VI sont quasiment absents des corpus d'entra\^inement
existants, alors que les phototypes~III--IV y sont d\'ej\`a partiellement
repr\'esent\'es). C'est pourquoi ce protocole concentre l'essentiel de
l'effort de collecte sur les phototypes~V et~VI plut\^ot que de r\'epartir
uniform\'ement l'effort sur l'ensemble de l'\'echelle.

\subsection{R\'epartition cible}

\begin{center}
\begin{tabularx}{\textwidth}{P{4cm} P{2.5cm} P{2.5cm} X}
\toprule
\textbf{Phototype} & \textbf{Part} & \textbf{Sujets} & \textbf{R\^ole} \\
\midrule
Fitzpatrick I--III & 15\% & 10--15 & Permet la comparaison directe avec la
litt\'erature existante (UBFC, PURE, VitalVideos-Europe-1) \\
Fitzpatrick IV & 20\% & 15--20 & Population de transition, teste la
continuit\'e des coefficients adapt\'es \\
Fitzpatrick V & 30\% & 20--25 & Priorit\'e haute --- cœur de la population
cible \\
Fitzpatrick VI & 35\% & 25--30 & Priorit\'e absolue --- population la plus
mal servie par les m\'ethodes existantes (MAE 14.1~bpm, Dasari et al.,
2021) \\
\bottomrule
\end{tabularx}
\end{center}

\subsection{Pourquoi ces seuils par groupe}

Pour qu'une comparaison statistique entre groupes de phototypes soit
significative, un minimum de 15 sujets par groupe est g\'en\'eralement
requis, 25 \'etant pr\'ef\'erable. Avec 70--100 sujets r\'epartis selon le
tableau ci-dessus, chaque groupe atteint ou d\'epasse ce seuil, y compris le
groupe le plus restreint (I--III).

\subsection{Classification du phototype --- m\'ethode \`a deux niveaux}

\subsubsection*{Niveau 1 : classification visuelle (sur le terrain)}

Un membre de l'\'equipe attribue une classe Fitzpatrick (1--6) au sujet au
moment de la collecte, par observation directe. Cette classification rapide
permet de suivre en temps r\'eel la progression vers les quotas du
tableau~4.3 et de prioriser le recrutement des phototypes manquants.

\subsubsection*{Niveau 2 : mesure objective (en post-traitement)}

La classification visuelle reste, par nature, un jugement subjectif --- une
limite que VitalVideos identifie explicitement dans son propre protocole
(Toye, 2023). Pour la compl\'eter, le vecteur de peau normalis\'e est
calcul\'e directement depuis les pixels d'une zone de joue homog\`ene, sur
une frame stable de l'enregistrement~1~:

\begin{codebox}
def mesurer\_vecteur\_peau(roi):\\
\quad\# roi~: tableau numpy [H, W, 3] en RGB\\
\quad R = roi[:,:,0].mean()\\
\quad G = roi[:,:,1].mean()\\
\quad B = roi[:,:,2].mean()\\[4pt]
\quad norme = sqrt(R**2 + G**2 + B**2)\\
\quad vecteur = [R, G, B] / norme\\[4pt]
\quad return vecteur\\
\quad \# Reference De Haan : [0.7682, 0.5121, 0.3841]
\end{codebox}

Une classification de phototype par seuils de luminosit\'e du canal vert
($G/255$) compl\`ete cette mesure~:

\begin{codebox}
luminosite = G / 255.0\\[2pt]
si luminosite \textgreater{} 0.50 : phototype = 1\\
sinon si luminosite \textgreater{} 0.42 : phototype = 2\\
sinon si luminosite \textgreater{} 0.35 : phototype = 3\\
sinon si luminosite \textgreater{} 0.28 : phototype = 4\\
sinon si luminosite \textgreater{} 0.22 : phototype = 5\\
sinon : phototype = 6
\end{codebox}

\begin{warnbox}{Les seuils ci-dessus sont indicatifs}
Ces seuils proviennent d'une estimation g\'en\'erale et doivent \^etre
recalibr\'es \`a partir des premi\`eres sessions du pilote (section~13), en
comparant la classification automatique \`a la classification visuelle
effectu\'ee sur le terrain. Le vecteur de peau mesur\'e (niveau~2) est de
toute mani\`ere la donn\'ee la plus utile pour le recalcul des coefficients
CHROM --- la classification en 6 classes n'est qu'un interm\'ediaire de
suivi des quotas.
\end{warnbox}

\newpage
\section{Mat\'eriel et r\'eglages}

\subsection{Smartphones --- rotation de 3 \`a 4 appareils}

\textbf{Choix recommand\'es}~: Samsung A03, Samsung A13, Tecno, Infinix ---
mod\`eles repr\'esentatifs du parc install\'e dans la r\'egion cibl\'ee.

\textbf{Pourquoi une rotation et pas un seul appareil.} Deux exemplaires
\guillemotleft~identiques\guillemotright\ d'un m\^eme mod\`ele peuvent avoir
des versions de firmware diff\'erentes selon la date et le lieu d'achat. Le
firmware contr\^ole des traitements internes --- balance des blancs
automatique, r\'eduction de bruit, niveau de compression --- qui peuvent
att\'enuer ou amplifier les micro-variations de couleur que le rPPG cherche
\`a capter. Utiliser plusieurs appareils en rotation, avec documentation
syst\'ematique du firmware, permet de mesurer cet effet a posteriori plut\^ot
que de le d\'ecouvrir trop tard sur un seul appareil non repr\'esentatif.

\textbf{Pour chaque appareil, documenter}~:
\begin{itemize}
\item Mod\`ele exact
\item Version Android et num\'ero de build (\emph{Param\`etres $\rightarrow$
\`A propos du t\'el\'ephone $\rightarrow$ Num\'ero de build})
\item Application cam\'era utilis\'ee
\item R\'esolution et bitrate vid\'eo configur\'es
\end{itemize}

\textbf{R\'eglages cam\'era recommand\'es.} R\'esolution et bitrate les plus
\'elev\'es disponibles dans l'application, en d\'esactivant toute option de
\guillemotleft~compression intelligente\guillemotright\ ou
\guillemotleft~optimisation du stockage\guillemotright. La compression
vid\'eo standard des smartphones (H.264) peut, \`a bitrate trop bas, lisser
des variations de pixel de l'ordre de 0.1\% --- exactement l'amplitude du
signal rPPG. \`A titre de r\'ef\'erence, un enregistrement de 30~secondes en
1080p devrait peser plusieurs dizaines \`a une centaine de m\'egaoctets~; un
fichier anormalement petit ($<$20~Mo pour 30~s) indique une compression
trop agressive.

\textbf{Distance et position.} Distance fixe de 50 \`a 70~cm entre le
smartphone et le visage, mat\'erialis\'ee par un rep\`ere visuel sur le
support ou la table. Hauteur de l'objectif au niveau des yeux du sujet.

\subsection{Oxym\`etre de pouls}

\textbf{Mod\`ele recommand\'e}~: type Contec CMS50D+ (connexion USB), comme
utilis\'e dans le protocole VitalVideos (Toye, 2023). \`A \'eviter~:
versions Bluetooth (latence variable, complique la synchronisation).

\textbf{Placement}~: au doigt du bras \emph{oppos\'e} \`a celui qui portera
le brassard de tension (voir section~5.3 et~9.2 pour la justification).

\subsection{Tensiom\`etre}

\textbf{Mod\`ele recommand\'e}~: tensiom\`etre \`a brassard de bras, type
Omron --- conforme au mat\'eriel utilis\'e par VitalVideos.

\textbf{Hygi\`ene}~: renouveler le brassard p\'eriodiquement. \`A l'\'echelle
de ce protocole (70--100 sujets), un renouvellement tous les 50 \`a 75 sujets
est raisonnable.

\subsection{Luxm\`etre}

Un luxm\`etre simple (capteur d\'edi\'e ou application smartphone calibr\'ee)
permet de mesurer objectivement l'\'eclairage ambiant au niveau du visage.

\textbf{Seuil de r\'ef\'erence}~: 150~lux minimum --- le m\^eme seuil que
celui utilis\'e par VitalVideos (Toye, 2023), qui ajoutait une lampe
annulaire lorsque la lumi\`ere disponible n'atteignait pas ce niveau.
En-dessous de ce seuil, en particulier pour les phototypes V--VI, le risque
de saturation basse des pixels (section~1.1) augmente significativement.

\textbf{Plage recommand\'ee pour phototypes V--VI}~: entre 200 et 800~lux.
Au-del\`a de 800~lux, la saturation haute devient un risque, notamment en
plein soleil direct (TYPE C).

\subsection{Indicateur de qualit\'e en temps r\'eel (recommand\'e)}

Pour \'eviter de d\'ecouvrir en post-traitement qu'un enregistrement est
inutilisable, une v\'erification simple peut \^etre effectu\'ee
imm\'ediatement apr\`es chaque enregistrement, sur la premi\`ere frame
stable~:

\begin{codebox}
def verifier\_qualite\_frame(frame\_rgb, face\_bbox):\\
\quad x, y, w, h = face\_bbox\\
\quad roi = frame\_rgb[y:y+h, x:x+w]\\[4pt]
\quad \# Proportion de pixels satures (\textgreater250 ou \textless5)\\
\quad satures = ((roi \textgreater\ 250) | (roi \textless\ 5)).mean()\\
\quad si satures \textgreater\ 0.15 :\\
\quad\quad retourner "Eclairage a corriger (trop fort ou trop faible)"\\[4pt]
\quad \# Taille du visage dans l'image (proxy de la distance)\\
\quad ratio\_visage = h / frame\_rgb.shape[0]\\
\quad si ratio\_visage \textless\ 0.15 :\\
\quad\quad retourner "Trop loin -- rapprocher le smartphone"\\
\quad si ratio\_visage \textgreater\ 0.85 :\\
\quad\quad retourner "Trop pres -- eloigner le smartphone"\\[4pt]
\quad retourner "Conditions acceptables"
\end{codebox}

Cette v\'erification permet de refaire imm\'ediatement un enregistrement
d\'efectueux plut\^ot que de le d\'ecouvrir lors de l'analyse, plusieurs
jours ou semaines plus tard.

\subsection{Support / tr\'epied l\'eger}

Un support stabilise le smartphone \`a distance constante et r\'eduit le
mouvement parasite, ce qui b\'en\'eficie \`a la fois au signal rPPG (moins
de glissement de la zone d'int\'er\^et) et au signal rBCG (le mouvement
cardiaque, de l'ordre de 0.01 \`a 0.1~mm, est noy\'e par tout mouvement de
plus grande amplitude --- voir section~8.3). Pour le TYPE C, le smartphone
peut \^etre tenu \`a la main pour repr\'esenter fid\`element l'usage r\'eel,
mais cela doit \^etre document\'e (champ \texttt{tenu\_main}/\texttt{pose}
dans les m\'etadonn\'ees).

\newpage
\section{Couvrir la diversit\'e de la lumi\`ere naturelle (TYPE C)}

\subsection{Pourquoi la lumi\`ere naturelle est trait\'ee \`a part}

Une analyse comparative sur le dataset MMPD montre que, pour les sources de
lumi\`ere artificielle (LED, fluorescent, incandescent), les diff\'erences
de performance entre m\'ethodes sont faibles --- mais que \emph{tous} les
mod\`eles performent significativement moins bien sous lumi\`ere naturelle.

\textbf{Cause principale identifi\'ee.} Aucune vid\'eo des datasets
d'entra\^inement de r\'ef\'erence (UBFC, AFRL) n'a \'et\'e enregistr\'ee sous
\'eclairage incandescent ou lumi\`ere naturelle~; leur composition spectrale
diff\`ere significativement de celle de la lumi\`ere du soleil. Les
mod\`eles n'ont donc jamais appris \`a reconna\^itre le
\guillemotleft~bruit\guillemotright\ caract\'eristique de ce type de
lumi\`ere, contrairement au bruit LED/fluorescent qu'ils savent d\'ej\`a
filtrer. Le m\^eme constat s'applique \`a VitalVideos-Europe-1, collect\'e
exclusivement en int\'erieur faute de conditions m\'et\'eorologiques
favorables (Toye, 2023) --- ce protocole corrige explicitement cette
limite (section~3.6).

\textbf{Implication directe.} Ce n'est pas une limite physique
incontournable --- une m\'ethode de personnalisation (MobilePhys) a
montr\'e une r\'eduction d'environ 40\% de l'erreur sous \'eclairage
incandescent simplement en exposant le mod\`ele \`a quelques exemples de
cette condition. La collecte TYPE C, en couvrant explicitement la lumi\`ere
naturelle, doit permettre un gain similaire par fine-tuning.

\subsection{Une seule journ\'ee de collecte ne suffit pas}

Le spectre de la lumi\`ere naturelle change continuellement~: angle du
soleil, pr\'esence de nuages, r\'eflexions ambiantes. Concentrer toute la
collecte TYPE C sur une seule journ\'ee, \`a une seule heure, ne couvrirait
qu'une combinaison parmi beaucoup d'autres --- l'\'equivalent de
n'enregistrer \guillemotleft~une seule LED\guillemotright\ en pensant
couvrir \guillemotleft~toutes les LED\guillemotright.

\subsection{Grille de r\'epartition recommand\'ee}

\begin{center}
\begin{tabularx}{\textwidth}{P{3.5cm} P{3cm} X}
\toprule
\textbf{Dimension} & \textbf{Valeurs} & \textbf{Justification} \\
\midrule
Moment de la journ\'ee & Matin (8h--10h), midi (12h--14h), fin de journ\'ee
(16h--18h) & La temp\'erature de couleur de la lumi\`ere solaire varie
fortement (plus chaude/orang\'ee en d\'ebut et fin de journ\'ee, plus
neutre/bleut\'ee \`a midi) \\
M\'et\'eo & Ensoleill\'e, nuageux & Le ciel nuageux diffuse la lumi\`ere
(spectre plus uniforme mais intensit\'e plus faible)~; le ciel d\'egag\'e
ajoute une forte composante directionnelle \\
Position relative au soleil & Ombre, plein soleil, lumi\`ere lat\'erale
privil\'egi\'ee & Le contre-jour direct sature le fond et sous-expose le
visage~; la lumi\`ere lat\'erale est la condition la plus repr\'esentative
d'un usage r\'eel \`a l'ext\'erieur d'une pharmacie \\
\bottomrule
\end{tabularx}
\end{center}

\subsection{Le contre-jour --- pi\`ege \`a \'eviter syst\'ematiquement}

\begin{warnbox}{R\`egle pratique}
Si le sujet a le soleil dans le dos, son visage est dans l'ombre alors que
le fond de l'image est tr\`es lumineux --- les pixels du fond saturent et le
visage est sous-expos\'e. Si le sujet fait face au soleil, l'\'eblouissement
provoque un plissement des yeux et des ombres dures (nez, sourcils) qui
augmentent le contraste local de mani\`ere non repr\'esentative.

\textbf{Consigne}~: orienter syst\'ematiquement le sujet avec la source de
lumi\`ere naturelle sur le c\^ot\'e (lumi\`ere lat\'erale). C'est \'egalement
la condition la plus probable lorsqu'un patient sort naturellement d'une
pharmacie en plein jour.
\end{warnbox}

\subsection{Avec seulement 5--6 sujets par cr\'eneau}

Avec 20--30 sujets TYPE C r\'epartis sur les trois moments de la journ\'ee et
les deux conditions m\'et\'eo, chaque combinaison re\c{c}oit environ
3--6~sujets. Ce n'est pas suffisant pour une analyse statistique fine par
combinaison, mais c'est suffisant pour que l'ensemble TYPE~C, dans son
ensemble, couvre une vraie diversit\'e spectrale --- ce qui est l'objectif
du fine-tuning (le mod\`ele voit des exemples vari\'es, m\^eme s'ils ne sont
pas nombreux individuellement).

\newpage
\section{De la collecte aux coefficients CHROM adapt\'es}

Cette section explique en d\'etail comment les donn\'ees collect\'ees
permettront de recalculer des coefficients CHROM adapt\'es par phototype ---
l'une des contributions scientifiques attendues de ce dataset.

\subsection{Rappel du pipeline CHROM standard}

\begin{enumerate}
\item \textbf{Normalisation}~: $R_n = R/\mu(R)$, $G_n = G/\mu(G)$,
$B_n = B/\mu(B)$ --- chaque canal oscille autour de 1. Cette \'etape
\'elimine le niveau moyen (DC) de la lumi\`ere, mais pas son oscillation
(AC)~: une lumi\`ere qui varie de mani\`ere p\'eriodique pendant
l'enregistrement reste pr\'esente apr\`es normalisation, \`a une \'echelle
diff\'erente.
\item \textbf{Standardisation par le vecteur de peau}~: $R_s = v_R R_n$,
$G_s = v_G G_n$, $B_s = v_B B_n$, avec $[v_R, v_G, v_B]$ le vecteur de peau
normalis\'e (standard~: $[0.7682, 0.5121, 0.3841]$).
\item \textbf{Projections de chrominance}~:
\[
X_s = \frac{R_s - G_s}{v_R - v_G}, \qquad
Y_s = \frac{R_s + G_s - 2B_s}{v_R + v_G - 2v_B}
\]
\item \textbf{Filtrage passe-bande} (0.75--3.5~Hz) sur $X_s$ et $Y_s$ pour
obtenir $X_f$ et $Y_f$.
\item \textbf{Combinaison adaptative}~: $S = X_f - \alpha Y_f$, avec
$\alpha = \sigma(X_f)/\sigma(Y_f)$.
\end{enumerate}

\subsection{D'o\`u viennent les coefficients $3, 2, 1.5, 1, 1.5$}

En d\'eveloppant les expressions de $X_s$ et $Y_s$ avec le vecteur standard
$[0.7682, 0.5121, 0.3841]$~:
\[
X_s = \frac{0.7682}{0.7682-0.5121}R_n - \frac{0.5121}{0.7682-0.5121}G_n
    = 3.0\,R_n - 2.0\,G_n
\]
\[
Y_s = \frac{0.7682}{0.5121}R_n + \frac{0.5121}{0.5121}G_n -
\frac{2\times0.3841}{0.5121}B_n
    = 1.5\,R_n + 1.0\,G_n - 1.5\,B_n
\]
Les coefficients $3, 2, 1.5, 1, 1.5$ sont donc enti\`erement d\'etermin\'es
par le choix du vecteur $[v_R, v_G, v_B]$. \textbf{Si ce vecteur change, les
coefficients changent selon la m\^eme formule.}

\subsection{La fonction de recalcul, \`a appliquer aux donn\'ees
collect\'ees}

Pour chaque sujet (ou chaque groupe de phototype, apr\`es agr\'egation des
vecteurs mesur\'es en section~4.6), de nouveaux coefficients sont
calcul\'es~:

\begin{codebox}
def calculer\_coefficients\_chrom(vecteur):\\
\quad vR, vG, vB = vecteur\\[4pt]
\quad denom\_X = vR - vG\\
\quad a1 = vR / denom\_X\quad\#\ remplace 3\\
\quad a2 = vG / denom\_X\quad\#\ remplace 2\\[4pt]
\quad denom\_Y = vR + vG - 2*vB\\
\quad a3 = vR / denom\_Y\quad\#\ remplace 1.5\\
\quad a4 = vG / denom\_Y\quad\#\ remplace 1.0\\
\quad a5 = 2*vB / denom\_Y\ \#\ remplace 1.5\\[4pt]
\quad return a1, a2, a3, a4, a5\\[6pt]
\#\ Verification : le vecteur De Haan original redonne\\
\#\ exactement [3.0, 2.0, 1.5, 1.0, 1.5]
\end{codebox}

\subsection{Ce que les donn\'ees collect\'ees doivent permettre de
produire}

\`A l'issue de la collecte, pour chaque groupe de phototype (I--III, IV, V,
VI), un tableau de coefficients recalcul\'es du type suivant doit pouvoir
\^etre produit~:

\begin{center}
\begin{tabularx}{\textwidth}{P{2.5cm} P{1.4cm} P{1.4cm} P{1.4cm} P{1.4cm}
P{1.4cm}}
\toprule
\textbf{Phototype} & $a_1$ & $a_2$ & $a_3$ & $a_4$ & $a_5$ \\
\midrule
I--II (r\'ef\'erence) & 3.00 & 2.00 & 1.50 & 1.00 & 1.50 \\
III & ? & ? & ? & ? & ? \\
IV & ? & ? & ? & ? & ? \\
V & ? & ? & ? & ? & ? \\
VI & ? & ? & ? & ? & ? \\
\bottomrule
\end{tabularx}
\end{center}

\textbf{Pour produire ce tableau, chaque sujet doit fournir}~:
\begin{itemize}
\item Un vecteur de peau mesur\'e dans des conditions d'\'eclairage
document\'ees et stables (id\'ealement TYPE A, o\`u l'\'eclairage est
contr\^ol\'e et le vecteur mesur\'e n'est pas confondu avec un effet de la
lumi\`ere ambiante).
\item Sa classe de phototype (section~4.5).
\end{itemize}

\subsection{Au-del\`a de la formule analytique : coefficients appris}

Une fois les coefficients analytiques recalcul\'es (7.4) sur le TYPE A, une
seconde \'etape consiste \`a les affiner par apprentissage --- en utilisant
les coefficients $a_1, \ldots, a_5$ comme param\`etres entra\^inables d'un
mod\`ele PyTorch, initialis\'es aux valeurs de De Haan, puis optimis\'es pour
minimiser l'\'ecart entre le HR pr\'edit et le HR de r\'ef\'erence (oxym\`etre)
sur les donn\'ees TYPE~A et TYPE~B du m\^eme groupe de phototype. Cette
\'etape ne n\'ecessite aucune donn\'ee suppl\'ementaire par rapport \`a ce
que ce protocole collecte d\'ej\`a --- elle r\'eutilise directement les
paires (signal vid\'eo, HR oxym\`etre) de chaque sujet.

\begin{notebox}{Pourquoi cette double approche est une contribution
originale}
\`A notre connaissance, aucune publication n'a mesur\'e le vecteur de peau
r\'eel ni recalcul\'e les coefficients CHROM pour des populations africaines
en conditions de pharmacie. La combinaison d'un recalcul analytique
(section~7.4, interpr\'etable, li\'e directement \`a la physique de la
m\'elanine) et d'un raffinement appris (section~7.5, qui peut capturer des
effets que la formule analytique ne mod\'elise pas, comme l'interaction avec
l'\'eclairage LED) constitue un r\'esultat scientifique exploitable
ind\'ependamment du fine-tuning des mod\`eles de deep learning.
\end{notebox}

\newpage
\section{Capture optionnelle du rBCG et estimation du PTT}

\subsection{Pourquoi consid\'erer le rBCG en compl\'ement du rPPG}

Le rPPG (variation de couleur) et le rBCG (remote
ballistocardiography --- micro-mouvement m\'ecanique caus\'e par
l'\'ejection du sang \`a chaque battement) sont deux signaux du m\^eme
\'ev\'enement cardiaque, mais sensibles \`a des sources de bruit
diff\'erentes~:

\begin{center}
\begin{tabularx}{\textwidth}{P{4cm} X X}
\toprule
\textbf{Condition} & \textbf{rPPG} & \textbf{rBCG} \\
\midrule
\'Eclairage variable & D\'egrad\'e --- la variation de lumi\`ere se confond
avec le signal de couleur & Non affect\'e --- le mouvement ne d\'epend pas de
la couleur \\
Peau tr\`es fonc\'ee & SNR r\'eduit --- la m\'elanine absorbe le signal
diffus & Non affect\'e --- la m\'elanine ne change pas la position des
pixels \\
Sujet immobile & Non affect\'e & Signal faible --- le micro-mouvement
cardiaque (0.01--0.1~mm) approche la r\'esolution du capteur \\
Sujet bouge (parole, rotation t\^ete) & Partiellement compens\'e par CHROM
& Masqu\'e --- un mouvement de quelques mm domine totalement le
micro-mouvement cardiaque \\
\bottomrule
\end{tabularx}
\end{center}

La capture simultan\'ee des deux signaux permettrait, en TYPE A et B, une
fusion pond\'er\'ee par le rapport signal/bruit de chacun, potentiellement
plus robuste qu'un seul des deux signaux pris isol\'ement --- en particulier
pour les phototypes V--VI o\`u le rPPG seul est le plus pénalis\'e.

\subsection{Le Pulse Transit Time (PTT) et la pression art\'erielle}

Le PTT est le d\'elai entre deux signaux li\'es au m\^eme battement
cardiaque, mesur\'es \`a deux endroits diff\'erents. Ce d\'elai est
corr\'el\'e \`a la vitesse de propagation de l'onde de pouls dans les
art\`eres, elle-m\^eme li\'ee \`a la rigidit\'e art\'erielle et donc \`a la
pression art\'erielle~:

\begin{center}
\begin{tabularx}{\textwidth}{P{2.8cm} P{2.8cm} P{2.8cm} X}
\toprule
\textbf{\'Etat} & \textbf{Vaisseaux} & \textbf{PTT} & \textbf{Pression} \\
\midrule
Hypertension & Rigides & Court ($<$100~ms) & \'Elev\'ee \\
Normo-tension & Normaux & Normal (100--200~ms) & Normale \\
Hypotension & Souples & Long ($>$200~ms) & Basse \\
\bottomrule
\end{tabularx}
\end{center}

Dans ce protocole, le PTT serait approxim\'e par le d\'elai entre un pic du
signal rBCG (visage) et le pic correspondant du signal rPPG, ou entre deux
signaux rPPG mesur\'es sur deux r\'egions diff\'erentes (visage et paume ---
d'o\`u l'int\'er\^et de la vid\'eo paume optionnelle de la section~2.2). Ce
PTT approxim\'e diff\`ere du PTT au sens strict (d\'elai ECG/PPG), comme
not\'e en section~2.3.

\begin{warnbox}{Statut de cette section}
La capture du rBCG et l'estimation du PTT restent \textbf{optionnelles} dans
ce protocole de collecte~: elles n'ajoutent aucun mat\'eriel suppl\'ementaire
(le rBCG est extrait de la m\^eme vid\'eo visage par optical flow), mais
elles ajoutent une charge de traitement et d'analyse. Si le temps disponible
par sujet est contraint, la priorit\'e reste les enregistrements~1 et~2
d\'ecrits en section~9.3--9.4. La vid\'eo paume (10--15~s suppl\'ementaires)
peut \^etre ajout\'ee en TYPE A et B si le temps le permet, sans
modification du mat\'eriel.
\end{warnbox}

\subsection{Limite \`a anticiper~: la stabilisation optique (OIS)}

Les smartphones haut de gamme corrigent activement les micro-mouvements de
la cam\'era (stabilisation optique, OIS), ce qui peut annuler pr\'ecis\'ement
le signal rBCG recherch\'e. Les smartphones d'entr\'ee de gamme
recommand\'es pour ce protocole (section~5.1) en sont g\'en\'eralement
d\'epourvus ou peu \'equip\'es --- ce qui, dans ce cas pr\'ecis, favorise la
collecte d'un signal rBCG exploitable.

\newpage
\section{Proc\'edure d\'etaill\'ee, \'etape par \'etape}

\subsection{\'Etape 1 --- Accueil et consentement}

\begin{enumerate}
\item Pr\'esentation orale du projet, en fran\c{c}ais et dans la langue
locale si n\'ecessaire~: objectif (am\'eliorer la mesure de signes vitaux par
smartphone pour les populations africaines), dur\'ee (environ 10--15~min),
caract\`ere anonyme et volontaire.
\item Signature du formulaire de consentement \'eclair\'e, incluant
explicitement~:
  \begin{itemize}
  \item l'usage des donn\'ees \`a des fins de recherche \textbf{et} de
  d\'eveloppement commercial de technologies de sant\'e num\'erique, dans le
  respect de l'anonymat complet des participants~;
  \item la possibilit\'e de partage des donn\'ees anonymis\'ees avec des
  partenaires de recherche~;
  \item le droit de retrait \`a tout moment, sans justification.
  \end{itemize}
\item Attribution d'un identifiant anonyme (ex.\ \texttt{CIV\_001}). Aucun
nom, num\'ero de t\'el\'ephone ou autre identifiant direct n'appara\^it dans
les fichiers de donn\'ees.
\end{enumerate}

\subsection{\'Etape 2 --- Installation}

\begin{enumerate}
\item Sujet assis, smartphone sur support \`a 50--70~cm, hauteur de
l'objectif au niveau des yeux (section~5.1).
\item Oxym\`etre plac\'e au doigt du \textbf{bras A} (sans brassard).
\item Brassard de tension plac\'e sur le \textbf{bras B}, \textbf{non
gonfl\'e} \`a ce stade.

\begin{warnbox}{Pourquoi cette s\'eparation des bras est imp\'erative}
Si l'oxym\`etre et le brassard sont sur le m\^eme bras, le gonflage du
brassard restreint temporairement la circulation sanguine dans ce bras et
provoque une perte du signal de l'oxym\`etre pendant et juste apr\`es la
mesure de pression. Cette r\`egle est directement reprise du protocole
VitalVideos (Toye, 2023). Placer les deux appareils sur des bras
diff\'erents \'evite cette interf\'erence et permet de capturer le PPG en
continu pendant les deux enregistrements.
\end{warnbox}

\item Mesure du lux ambiant au niveau du visage avec le luxm\`etre (cible
$\geq$150~lux, section~5.4).
\item Classification visuelle du phototype Fitzpatrick (1--6, section~4.5).
\item Si une vid\'eo paume est pr\'evue (section~8.2), pr\'eparer la
position de la main \`a proximit\'e de la zone de cadrage.
\end{enumerate}

\subsection{\'Etape 3 --- Enregistrement~1 : \'etat de repos
(60--90~secondes)}

\begin{enumerate}
\item V\'erifier que le brassard reste \textbf{non gonfl\'e} pendant toute
cette \'etape.
\item D\'emarrer simultan\'ement~:
  \begin{itemize}
  \item l'enregistrement vid\'eo du visage (smartphone)~;
  \item la capture PPG/HR/SpO2 (oxym\`etre).
  \end{itemize}
\item \textbf{Synchronisation}~: pendant les 3 premi\`eres secondes de la
vid\'eo, tenir l'oxym\`etre face \`a la cam\'era et d\'eclencher son
affichage. Ce moment, visible dans la vid\'eo, servira de rep\`ere temporel
commun ($t=0$) lors du post-traitement --- aucun mat\'eriel de
synchronisation suppl\'ementaire n'est n\'ecessaire (m\^eme principe que
VitalVideos, qui synchronise via la connexion USB directe de sa cam\'era~;
ici on substitue un rep\`ere visuel car la cam\'era est celle du
smartphone).
\item Consigne selon le type de collecte (section~3)~:
  \begin{itemize}
  \item \textbf{TYPE A}~: immobile, silencieux, regard fixe vers la
  cam\'era.
  \item \textbf{TYPE B}~: position assise naturelle, parole douce
  autoris\'ee.
  \item \textbf{TYPE C}~: conditions ext\'erieures ou r\'eelles, l\'egers
  mouvements naturels tol\'er\'es (rotation l\'eg\`ere de la t\^ete, parole).
  \end{itemize}
\item \`A la fin de l'enregistrement, v\'erifier imm\'ediatement la
qualit\'e (section~5.5) sur une frame de la vid\'eo. En cas de probl\`eme,
recommencer cet enregistrement avant de passer \`a l'\'etape suivante.
\end{enumerate}

\subsection{\'Etape 4 --- Enregistrement~2 : mesure de la pression
art\'erielle (30--60~secondes)}

\begin{enumerate}
\item D\'emarrer simultan\'ement~: l'enregistrement vid\'eo, le gonflage du
brassard, et la mesure de pression art\'erielle.
\item La capture PPG/HR/SpO2 continue pendant cette phase.

\begin{warnbox}{Pourquoi la pression doit \^etre mesur\'ee \emph{pendant}
l'enregistrement, jamais avant ni apr\`es}
Le gonflage du brassard est inconfortable et peut provoquer une l\'eg\`ere
augmentation transitoire de la fr\'equence cardiaque et de la pression ---
un effet de stress li\'e \`a la mesure elle-m\^eme. Si la pression est
mesur\'ee \emph{avant} l'enregistrement~1, ce dernier ne refl\`ete plus
l'\'etat de repos r\'eel. Si elle est mesur\'ee \emph{apr\`es}, la valeur de
pression ne correspond \`a aucune vid\'eo. VitalVideos (Toye, 2023) souligne
explicitement ce point comme \guillemotleft~de la plus haute
importance\guillemotright\ dans son protocole. En la mesurant
simultan\'ement avec l'enregistrement~2 --- un enregistrement distinct,
explicitement \'etiquet\'e comme correspondant \`a cet \'etat --- les deux
enregistrements restent coh\'erents avec l'\'etat physiologique qu'ils sont
suppos\'es repr\'esenter.
\end{warnbox}

\item Cette vid\'eo (enregistrement~2) sert sp\'ecifiquement \`a la
calibration de la pression art\'erielle pour le PTT (section~8.2)~; elle
n'est pas suppos\'ee repr\'esenter un \'etat de repos pur.
\end{enumerate}

\subsection{\'Etape 5 --- (Optionnelle) Vid\'eo paume}

Si pr\'evue (section~8.2)~: courte vid\'eo additionnelle de la paume de la
main, 10--15~secondes, m\^eme smartphone, m\^emes conditions d'\'eclairage
que les enregistrements pr\'ec\'edents.

\subsection{\'Etape 6 --- Cl\^oture}

\begin{enumerate}
\item Affichage au participant des valeurs mesur\'ees (HR, SpO2, pression
art\'erielle) --- transparence et retour direct sur la session, comme
pratiqu\'e par VitalVideos (Toye, 2023).
\item Remerciement, et compensation \'eventuelle selon les modalit\'es
d\'efinies avec la pharmacie.
\item V\'erification finale que tous les fichiers (vid\'eos, export
oxym\`etre, m\'etadonn\'ees) sont correctement enregistr\'es sous
l'identifiant anonyme du sujet avant de passer au participant suivant.
\end{enumerate}

\newpage
\section{M\'etadonn\'ees enregistr\'ees par session}

\subsection{Pourquoi documenter autant}

Chaque variable suppl\'ementaire enregistr\'ee est une dimension d'analyse
possible au moment du fine-tuning et de l'\'evaluation~: erreur par
phototype, par type d'\'eclairage, par appareil, par niveau de mouvement.
Sans ces m\'etadonn\'ees, un mod\`ele qui performe mal ne pourrait pas \^etre
diagnostiqu\'e --- s'agit-il d'un probl\`eme li\'e au phototype, \`a
l'\'eclairage, \`a l'appareil, ou \`a une combinaison~?

\subsection{Structure recommand\'ee (par session)}

\begin{codebox}
\{\\
\ \ "id\_anonyme": "CIV\_001",\\
\ \ "date": "2026-07-15",\\
\ \ "heure": "09:30",\\
\ \ "type\_collecte": "A | B | C",\\
\ \ "lieu": "interieur\_pharmacie | exterieur | espace\_controle",\\[2pt]
\ \ "eclairage": \{\\
\ \ \ \ "type": "LED\_neutre | LED\_chaude | neon | naturel | mixte",\\
\ \ \ \ "lux\_mesure": 320,\\
\ \ \ \ "stable": true,\\
\ \ \ \ "orientation": "face | laterale | contre\_jour",\\
\ \ \ \ "moment\_journee": "matin | midi | fin\_journee",\\
\ \ \ \ "meteo": "ensoleille | nuageux | interieur"\\
\ \ \},\\[2pt]
\ \ "smartphone": \{\\
\ \ \ \ "modele": "Samsung A13",\\
\ \ \ \ "android\_version": "12",\\
\ \ \ \ "build": "A135FXXU2BWA1",\\
\ \ \ \ "app\_camera": "Open Camera",\\
\ \ \ \ "resolution": "1920x1080",\\
\ \ \ \ "bitrate\_estime\_Mbps": 35,\\
\ \ \ \ "support": "trepied | tenu\_main"\\
\ \ \},\\[2pt]
\ \ "sujet": \{\\
\ \ \ \ "age\_tranche": "30-40",\\
\ \ \ \ "genre": "M | F",\\
\ \ \ \ "phototype\_visuel": 6,\\
\ \ \ \ "vecteur\_peau\_calcule": [0.85, 0.43, 0.30],\\
\ \ \ \ "pathologies\_declarees": ["hypertension"],\\
\ \ \ \ "medicaments": []\\
\ \ \},\\[2pt]
\ \ "scenario": \{\\
\ \ \ \ "position": "assis | debout",\\
\ \ \ \ "mouvement": "immobile | leger | parole",\\
\ \ \ \ "distance\_cm": 60\\
\ \ \},\\[2pt]
\ \ "enregistrement\_1": \{\\
\ \ \ \ "duree\_s": 75,\\
\ \ \ \ "hr\_oxymetre": 72,\\
\ \ \ \ "spo2": 98,\\
\ \ \ \ "qualite\_signal\_ref": "bonne | moyenne | artefacts",\\
\ \ \ \ "fichier\_video": "CIV\_001\_rec1.mp4",\\
\ \ \ \ "fichier\_ppg": "CIV\_001\_rec1\_ppg.csv"\\
\ \ \},\\[2pt]
\ \ "enregistrement\_2": \{\\
\ \ \ \ "duree\_s": 40,\\
\ \ \ \ "bp\_systolique": 128,\\
\ \ \ \ "bp\_diastolique": 82,\\
\ \ \ \ "hr\_oxymetre": 75,\\
\ \ \ \ "spo2": 98,\\
\ \ \ \ "fichier\_video": "CIV\_001\_rec2.mp4",\\
\ \ \ \ "fichier\_ppg": "CIV\_001\_rec2\_ppg.csv"\\
\ \ \},\\[2pt]
\ \ "video\_paume": null\\
\}
\end{codebox}

\subsection{Champ \texttt{qualite\_signal\_ref}}

Ce champ, particuli\`erement important pour le TYPE~C (section~3.5),
permet de signaler si le signal de l'oxym\`etre a pr\'esent\'e des
artefacts visibles (perte de signal, valeurs aberrantes) pendant
l'enregistrement. Lors de l'entra\^inement, ce champ peut \^etre utilis\'e
pour pond\'erer la confiance accord\'ee \`a chaque label, ou pour exclure
les sessions les plus probl\'ematiques d'un sous-ensemble d'\'evaluation
prioritaire.

\newpage
\section{Comment ces donn\'ees seront utilis\'ees pour l'entra\^inement}

\subsection{Progression A $\rightarrow$ B $\rightarrow$ C}

La strat\'egie d'entra\^inement suit directement la logique des trois types
de collecte~:

\begin{enumerate}
\item \textbf{Pr\'e-entra\^inement} sur des datasets publics existants
(UBFC, PURE) --- le mod\`ele acquiert une connaissance g\'en\'erale du
probl\`eme rPPG.
\item \textbf{Fine-tuning sur TYPE~A} --- les labels les plus pr\'ecis
permettent un apprentissage fondamental propre, sp\'ecifique aux phototypes
V--VI.
\item \textbf{Fine-tuning sur TYPE~A~+~B} --- le mod\`ele s'adapte \`a la
variabilit\'e d'\'eclairage r\'eelle de la pharmacie, avec un taux
d'apprentissage r\'eduit pour ne pas perdre l'acquis de l'\'etape
pr\'ec\'edente.
\item \textbf{Fine-tuning final sur TYPE~A~+~B~+~C} --- dernier ajustement
pour les conditions exactes de d\'eploiement, avec un taux d'apprentissage
encore plus faible.
\end{enumerate}

\begin{warnbox}{Pourquoi pas un entra\^inement s\'equentiel strict
(A puis B puis C s\'epar\'ement)}
Un entra\^inement purement s\'equentiel --- d'abord uniquement A, puis
uniquement B, puis uniquement C --- risque un oubli catastrophique~: le
mod\`ele peut perdre la performance acquise sur les conditions pr\'ec\'edentes
en s'adaptant aux nouvelles. La recommandation est de \textbf{m\'elanger}
progressivement les ensembles (A, puis A+B, puis A+B+C) plut\^ot que de les
remplacer.
\end{warnbox}

\subsection{R\`egle d'\'evaluation}

Les donn\'ees d'\'evaluation ne doivent jamais avoir \'et\'e vues pendant
l'entra\^inement, et l'\'evaluation finale doit porter en priorit\'e sur le
sous-ensemble TYPE~C --- la condition la plus proche du d\'eploiement r\'eel.

\begin{center}
\begin{tabularx}{\textwidth}{P{3cm} P{3cm} X}
\toprule
\textbf{Type} & \textbf{Split recommand\'e} & \textbf{Justification} \\
\midrule
TYPE A & 80\% train / 20\% test & Fondation~; un petit ensemble de test
suffit pour v\'erifier l'absence de r\'egression \\
TYPE B & 70\% train / 30\% test & Cœur du dataset~; \'echantillon de test
plus large pour une mesure stable \\
TYPE C & 60\% train / 40\% test & Cible du d\'eploiement~; la plus grande
proportion de test pour une mesure fiable de la performance r\'eelle \\
\bottomrule
\end{tabularx}
\end{center}

\subsection{M\'etriques \`a calculer, et pourquoi chacune}

\begin{itemize}
\item \textbf{MAE globale} --- erreur absolue moyenne en bpm, comparable
directement aux chiffres de la litt\'erature (Dasari et al., 2021~: 5.2~bpm
pour CHROM sur phototypes~I--III, 14.1~bpm sur~V--VI).
\item \textbf{MAE par groupe de phototype} --- v\'erifie que
l'am\'elioration cibl\'ee sur V--VI ne se fait pas au d\'etriment des autres
groupes.
\item \textbf{MAE par condition d'\'eclairage} (LED neutre, LED chaude,
n\'eon, naturel) --- identifie si une condition particuli\`ere reste
probl\'ematique malgr\'e le fine-tuning, ce qui orienterait une couche de
contr\^ole qualit\'e en amont du mod\`ele plut\^ot qu'un nouvel
entra\^inement.
\item \textbf{MAE par niveau de mouvement} (immobile, l\'eger, parole) ---
mesure la robustesse au sc\'enario r\'eel d'utilisation.
\item \textbf{RMSE en compl\'ement de la MAE} --- un \'ecart important entre
RMSE et MAE signale la pr\'esence de quelques erreurs tr\`es importantes
(potentiellement des sessions \`a \texttt{qualite\_signal\_ref} d\'egrad\'ee,
section~10.3) plut\^ot qu'une erreur uniform\'ement r\'epartie.
\end{itemize}

\subsection{Tableau de comparaison final attendu}

\begin{center}
\begin{tabularx}{\textwidth}{X P{2.6cm} P{2.6cm} P{2.4cm}}
\toprule
\textbf{M\'ethode} & MAE peaux claires & MAE peaux fonc\'ees &
D\'eployable smartphone \\
\midrule
CHROM De Haan original & $\sim$5~bpm & $\sim$14~bpm & Oui \\
CHROM adapt\'e par phototype (section~7) & $\sim$4--5~bpm & objectif
$<$9~bpm & Oui \\
TS-CAN pr\'e-entra\^in\'e (sans fine-tuning) & $\sim$2~bpm & $\sim$10~bpm &
Oui \\
TS-CAN fine-tun\'e (A+B+C) & $\sim$2--3~bpm & objectif $<$5~bpm & Oui \\
\bottomrule
\end{tabularx}
\end{center}

\newpage
\section{Points de vigilance --- synth\`ese}

\begin{warnbox}{R\`egle absolue}
Ne jamais mesurer la pression art\'erielle avant ou apr\`es l'enregistrement
vid\'eo --- uniquement pendant l'enregistrement~2 (section~9.4), pour
\'eviter que l'inconfort du gonflage du brassard ne biaise l'\'etat de repos
de l'enregistrement~1.
\end{warnbox}

\begin{itemize}
\item L'oxym\`etre doit imp\'erativement \^etre plac\'e sur le bras
\textbf{sans} le brassard de tension (section~9.2).
\item Documenter pr\'ecis\'ement le mod\`ele et le firmware du smartphone
pour chaque session (section~5.1) --- permet d'analyser \emph{a posteriori}
l'impact de l'appareil sur la qualit\'e du signal.
\item Pour le TYPE~C, privil\'egier syst\'ematiquement la lumi\`ere
lat\'erale et \'eviter le contre-jour direct (section~6.4).
\item V\'erifier la qualit\'e de chaque enregistrement imm\'ediatement
apr\`es sa capture (section~5.5) --- recommencer plut\^ot que de
d\'ecouvrir un probl\`eme en post-traitement.
\item Maintenir l'\'eclairage entre 150 et 800~lux pour les phototypes
V--VI, en particulier en TYPE~A et~B (section~5.4) --- en-dessous, risque de
saturation basse~; au-dessus, risque de saturation haute.
\item M\'elanger les TYPE~A/B/C dans le fine-tuning final plut\^ot que de
les traiter s\'equentiellement (section~11.1), pour \'eviter l'oubli
catastrophique.
\item V\'erifier r\'eguli\`erement, pendant la collecte, la progression vers
les quotas de phototype (section~4.3) --- prioriser le recrutement des
groupes sous-repr\'esent\'es plut\^ot que de les rattraper en fin de
collecte.
\end{itemize}

\newpage

\section{Pilote recommand\'e avant d\'eploiement large}

\subsection{Objectif du pilote}

Avant la collecte \`a grande \'echelle, ce protocole complet doit \^etre
test\'e sur 3 \`a 5~sujets (coll\`egues ou volontaires), en suivant
exactement les \'etapes de la section~9, y compris la saisie compl\`ete des
m\'etadonn\'ees (section~10).

\subsection{Ce que le pilote doit permettre de v\'erifier}

\begin{enumerate}
\item \textbf{Dur\'ee r\'eelle par session} --- la cible indicative de
10--15~minutes par sujet (accueil, installation, deux enregistrements,
cl\^oture) doit \^etre confirm\'ee ou ajust\'ee.
\item \textbf{D\'elai de stabilisation de l'oxym\`etre} --- typiquement
10--15~secondes apr\`es la pose sur le doigt~; v\'erifier que ce d\'elai est
bien int\'egr\'e avant le d\'ebut effectif de l'enregistrement~1.
\item \textbf{Classification automatique du phototype} (section~4.5) ---
comparer les r\'esultats de \texttt{mesurer\_vecteur\_peau} et de la
classification par seuils \`a la classification visuelle effectu\'ee sur le
terrain~; ajuster les seuils si n\'ecessaire.
\item \textbf{Probl\`emes pratiques de d\'etection du visage} --- lunettes
(reflets), occlusions partielles, angle de cam\'era~; identifier les cas qui
posent probl\`eme \`a la d\'etection automatique avant la collecte \`a grande
\'echelle.
\item \textbf{Indicateur de qualit\'e en temps r\'eel} (section~5.5) ---
v\'erifier que les seuils de saturation et de distance propos\'es
correspondent \`a des probl\`emes r\'eellement visibles sur les vid\'eos.
\item \textbf{Confort du protocole de synchronisation} (section~9.3) ---
v\'erifier que le geste de pr\'esenter l'oxym\`etre face \`a la cam\'era
pendant 3~secondes est facile \`a r\'ealiser et bien visible en
post-traitement.
\end{enumerate}

\subsection{Ajustements attendus apr\`es le pilote}

Ce document constitue la \textbf{version~1} du protocole. Les param\`etres
suivants sont les plus susceptibles d'\^etre ajust\'es \`a la suite du
pilote~:

\begin{itemize}
\item Dur\'ees exactes des enregistrements~1 et~2 (actuellement 60--90~s et
30--60~s).
\item Seuils de classification automatique du phototype (section~4.5).
\item Seuils de l'indicateur de qualit\'e temps r\'eel (section~5.5).
\item Liste finale des smartphones en rotation, selon leur disponibilit\'e
r\'eelle et la qualit\'e du signal observ\'ee (section~5.1).
\end{itemize}

\newpage

\section*{R\'ef\'erences}
\addcontentsline{toc}{section}{R\'ef\'erences}

\begin{itemize}
\item De Haan, G. \& Jeanne, V. (2013). Robust pulse rate from
chrominance-based rPPG. \emph{IEEE Transactions on Biomedical Engineering},
60(10), 2878--2886.
\item Dasari, A., Prakash, S. K. A., Jeni, L. A. \& Tucker, C. S. (2021).
Evaluation of biases in remote photoplethysmography methods. \emph{npj
Digital Medicine}, 4, 91.
\item Mehmood, A. et al. (2025). Demographic bias in public remote
photoplethysmography datasets. \emph{npj Digital Medicine}.
\item Toye, P. (2023). \emph{Vital Videos: A public dataset of videos with
PPG and BP ground truths}. Vital Videos project.
\item Liu, X. et al. (2020). Multi-task temporal shift attention networks
for on-device contactless vitals measurement (TS-CAN). \emph{NeurIPS}.
\item Yu, Z. et al. (2019). RhythmNet / PhysNet : remote photoplethysmography
via spatio-temporal networks. \emph{BMVC}.
\end{itemize}

\vspace{0.6cm}
\noindent\rule{\textwidth}{0.4pt}
\begin{center}
\small Fin du document --- Protocole version 1, sujet \`a r\'evision apr\`es
le pilote (section~13)
\end{center}

\end{document}